\documentclass[11pt]{article}

\usepackage[margin=1in]{geometry}
\usepackage{amsmath,amssymb,amsthm}
\usepackage[hidelinks]{hyperref}
\hypersetup{
  pdftitle={An Aligned-Center Recursion for Products of q-Integers},
  pdfauthor={hainzhao}
}

\newtheorem{theorem}{Theorem}
\newtheorem{lemma}[theorem]{Lemma}
\newcommand{\qint}[1]{[#1]_q}
\newcommand{\qrint}[2]{[#1]_{q^{#2}}}
\newcommand{\archivedoi}{10.5281/zenodo.21830407}

\title{An Aligned-Center Recursion for Products of $q$-Integers}
\author{hainzhao}
\date{August 2026\\[4pt]
\small Proof and source archive:
\href{https://doi.org/\archivedoi}{doi:\archivedoi}}

\begin{document}
\maketitle

\begin{abstract}
Connelly, Ito, Martinez, Shevchenko, and Yang conjectured a sufficient
condition for the unimodality of
\(\prod_{i=1}^k[a_i]_q[b]_{q^r}\), and proved necessity in two
low-dimensional regimes.  We prove their sufficient condition for every
\(k\geq1\) and \(r\geq2\).  The proof is an induction based on the identity
\[
 [a+r]_q[b+1]_{q^r}
 =q^r[a]_q[b]_{q^r}+[a+r(b+1)]_q.
\]
After multiplication by the remaining factors, the two terms on the right
are symmetric unimodal about exactly the same center.  Allocating the
available increments among the parameters completes the induction.
\end{abstract}

\section{Introduction}

For a positive integer \(n\), write
\[
 [n]_q=1+q+\cdots+q^{n-1}.
\]
Products of such polynomials arise naturally in questions about
\(q\)-Fibonomial coefficients.  Connelly, Ito, Martinez, Shevchenko, and
Yang formulated the following sufficient condition
\cite[Conjecture~5.4]{CIMSY}: if \(r\geq2\), \(k\geq1\), and
\(a_1,\ldots,a_k,b\) are positive integers satisfying
\[
 r\mid a_i\quad\text{for some }i,
 \qquad\text{or}\qquad
 b\leq1+\sum_{i=1}^k\left\lfloor\frac{a_i}{r}\right\rfloor,
 \tag{1}
\]
then
\[
 P(q)=\prod_{i=1}^k[a_i]_q[b]_{q^r}                 \tag{2}
\]
is unimodal.  They also conjecture necessity when \(k\leq3\) or
\(r\leq3\), prove several precursor cases, and verify bounded instances.

We prove the universally quantified sufficient statement in (1).  In
particular, this settles its first corner beyond both low-dimensional
regimes, \(k=r=4\), but the argument does not depend on either value.

\begin{theorem}\label{thm:main}
Let \(r\geq2\), \(k\geq1\), and let \(a_1,\ldots,a_k,b\) be positive
integers.  If either condition in \emph{(1)} holds, then the polynomial
\emph{(2)} is symmetric and unimodal.
\end{theorem}

The condition is not necessary for general \(k,r\).  For example, the
source observes that \(([3]_q)^4[2]_{q^4}\) is unimodal although it violates
the inequality in (1) \cite[after Conjecture~5.4]{CIMSY}.  Thus
Theorem~\ref{thm:main} proves only the sufficient direction and makes no new
necessity claim.

\section{Two elementary lemmas}

A finite nonnegative coefficient sequence is \emph{symmetric unimodal about
\(C\)} if coefficients at exponents equidistant from \(C\) agree and the
coefficients weakly increase up to \(C\).  We allow initial and terminal
zeros; this convention makes translates convenient.

\begin{lemma}\label{lem:closure}
The product of two nonnegative symmetric unimodal polynomials is symmetric
unimodal.  A sum of nonnegative symmetric unimodal polynomials having a
common center is symmetric unimodal.
\end{lemma}

\begin{proof}
The assertion about sums follows coefficientwise.  For products, decompose
each symmetric unimodal coefficient sequence into a nonnegative linear
combination of indicator sequences of nested intervals having the same
center.  The convolution of two interval indicators is a symmetric
trapezoidal sequence, hence unimodal, and all resulting convolutions have
the sum of the original centers.  Summing them proves the claim.
\end{proof}

The next identity supplies the induction.

\begin{lemma}[aligned-center recursion]\label{lem:recursion}
For positive integers \(a,b,r\),
\[
 [a+r]_q[b+1]_{q^r}
 =q^r[a]_q[b]_{q^r}+[a+r(b+1)]_q.                  \tag{3}
\]
\end{lemma}

\begin{proof}
Use
\[
 [a+r]_q=[r]_q+q^r[a]_q,
 \qquad
 [b+1]_{q^r}=1+q^r[b]_{q^r}.
\]
After subtracting the first term on the right side of (3), the remaining
terms are
\begin{align*}
 [r]_q[b+1]_{q^r}+q^r[a]_q
       \bigl(1+(q^r-1)[b]_{q^r}\bigr)
 &= [r(b+1)]_q+q^{r(b+1)}[a]_q\\
 &= [a+r(b+1)]_q.
\end{align*}

For a coefficient-level check independent of this manipulation, the left
side counts pairs \((x,j)\) with \(0\leq x<a+r\) and \(0\leq j\leq b\),
weighted by \(x+rj\).  The pairs with \(x<a\) and \(j\geq1\) give
\(q^r[a]_q[b]_{q^r}\).  The complementary pairs have either \(j=0\), or
\(x\geq a\) and \(j\geq1\); their weights occur once each and fill the
consecutive interval from \(0\) to \(a+r(b+1)-1\).  They give the other
term in (3).
\end{proof}

\section{Proof of the theorem}

Suppose first that \(r\mid a_i\) for some \(i\).  Writing \(a_i=rs\), we
have
\[
 [a_i]_q[b]_{q^r}
 =[r]_q[s]_{q^r}[b]_{q^r}.
\]
The product \([s]_z[b]_z\) is symmetric unimodal in \(z\).  Substitution
\(z=q^r\), followed by multiplication by \([r]_q\), repeats each of its
coefficients in a block of length \(r\); the resulting sequence remains
symmetric unimodal.  Lemma~\ref{lem:closure} then permits multiplication by
the remaining factors \([a_j]_q\).  This proves the divisibility branch.

Now assume no \(a_i\) is divisible by \(r\), and that the inequality in (1)
holds.  Put \(c_i=\lfloor a_i/r\rfloor\).  Choose integers \(d_i\) such
that
\[
 0\leq d_i\leq c_i,
 \qquad \sum_{i=1}^k d_i=b-1.                      \tag{4}
\]
Such a choice exists by (1).  Define \(a_i^{(0)}=a_i-rd_i\).  Every
\(a_i^{(0)}\) is positive because \(r\nmid a_i\).  At spacer length one,
the base polynomial is
\[
 \prod_{i=1}^k[a_i^{(0)}]_q,
\]
which is symmetric unimodal by Lemma~\ref{lem:closure}.

Starting from this base, perform the \(b-1\) increments allocated in (4).
At one step, let \(a\) be the current value of the selected length, let
\(c\) be the current spacer length, and let \(B(q)\) be the product of all
other current ordinary \(q\)-integers.  Lemma~\ref{lem:recursion} gives
\[
 [a+r]_q B(q)[c+1]_{q^r}
 =q^r[a]_qB(q)[c]_{q^r}
  +[a+r(c+1)]_qB(q).                               \tag{5}
\]
By induction, the first term on the right is a translate of a symmetric
unimodal polynomial.  The second term is a product of ordinary
\(q\)-integers, hence is symmetric unimodal.

It remains only to check their centers.  If the old degree is
\[
 E=(a-1)+\deg B+r(c-1),
\]
then the new degree is \(E+2r\).  The support endpoints of the translated
first term sum to \(E+2r\), while the degree of the second term is
\[
 a+r(c+1)-1+\deg B=E+2r.
\]
Thus both summands in (5) have center \((E+2r)/2\).  Their sum is symmetric
unimodal by Lemma~\ref{lem:closure}.  Iterating (5) reconstructs all final
lengths \(a_i\) and the final spacer length \(b\), proving
Theorem~\ref{thm:main}.

\section{Reproducibility and scope}

The identity and induction use no computation.  The accompanying exact
standard-library replay nevertheless checks (3) by both polynomial
multiplication and the independent pair partition, checks center alignment
at every recursive step, and compares the recursive construction with direct
multiplication over a bounded regression corpus.  From the archive root run
\begin{verbatim}
python3 proof/qanalog_conjecture54_sufficiency.py
\end{verbatim}
The bounded corpus is regression evidence only; it is not an assumption in
the universal proof.

The theorem proves the sufficient direction of Conjecture~5.4 for all
\(k,r\).  It does not prove the conjectured necessity for \(k\leq3\) or
\(r\leq3\), does not characterize the additional unimodal products outside
(1), and does not by itself settle the general \(q\)-Fibonomial unimodality
conjecture.

\paragraph{AI assistance.}
The author used AI assistance in developing the proof presentation and its
verification code.

\bibliographystyle{plain}
\bibliography{references}

\end{document}
